\section{Infrastructure Overview}

To maintain the "Low Cost" objective during deployment, Origin AI utilizes a modern, containerized microservice architecture. This allows for lightweight deployment on commodity hardware or spot cloud instances.

\section{Containerization Strategy}

Both the \textbf{Next.js Frontend} and the \textbf{Python AI Microservice} are Dockerized to ensure environment parity across development, staging, and production.

\begin{itemize}
    \item \textbf{Next.js Frontend:} A multi-stage Docker build is used to keep the production image under 200MB. We use a Node.js \texttt{alpine} base to minimize the security attack surface.
    \item \textbf{Python Backend:} The microservice is packaged with FastAPI and LangChain. We use a \texttt{python:3.10-slim} base for optimal performance and smaller image size.
\end{itemize}

\section{Docker Compose Orchestration}

For local development and testing, we use \textbf{Docker Compose} to manage the lifecycle of the services:
\begin{enumerate}
    \item \texttt{origin-frontend}: The Next.js interaction layer.
    \item \texttt{origin-backend}: The Python orchestrator.
    \item \texttt{origin-db}: PostgreSQL with the \texttt{pgvector} extension.
\end{enumerate}

\section{Cloud Hosting and Scaling}

To handle sudden spikes in student traffic (e.g., during exam seasons), the infrastructure is designed for \textbf{Horizontal Scaling}.

\begin{itemize}
    \item \textbf{Statelessness:} The Python orchestrator is entirely stateless. All session memory is persisted in PostgreSQL, allowing any replica to pick up a conversation given the session ID.
    \item \textbf{Spot Instances:} Since AI generation is compute-heavy, we run backend containers on cloud spot instances, which are significantly cheaper than on-demand instances.
    \item \textbf{Serverless Fallback:} For low-traffic subjects, the backend can be deployed as a serverless function (e.g., Google Cloud Run) to ensure we pay zero during idle hours.
\end{itemize}
